% Additional faithful author context for the main plane tester, official Sections2–3 and Corollary52.
\section*{Additional author algebraic context}
Recall The definition of the binomial and multinomial coefficients for natural numbers:
\[\binom ab\stackrel{\mathrm{def}}=\frac{a!}{b!(a-b)!},\qquad\binom a{b_1,\ldots,b_\ell}=\frac{a!}{b_1!\cdots b_\ell!},\]
where $\sum b_i=a$. We extend this definition to $\boldsymbol I,\boldsymbol I_1,\boldsymbol J,\boldsymbol J_1,\ldots,\boldsymbol J_\ell\in\mathbb N^m$ where $\boldsymbol J=\boldsymbol J_1+\ldots+\boldsymbol J_\ell$ by
\[\binom{\boldsymbol I}{\boldsymbol I_1}\stackrel{\mathrm{def}}=\prod_{k=1}^m\binom{I_k}{(I_1)_k},\qquad\binom{\boldsymbol J}{\boldsymbol J_1,\ldots,\boldsymbol J_\ell}\stackrel{\mathrm{def}}=\prod_{k=1}^m\binom{J_k}{(J_1)_k,\cdots,(J_\ell)_k}.\]
We mention Lucas theorem, that if $p$ is prime, $q=p^w$ for $w\in\mathbb N$, and $a,b\in\mathbb N$ have base $p$ representation $a=\sum_{\ell=0}^wa_\ell p^\ell$, $b=\sum_{\ell=0}^wb_\ell p^\ell$ with $0\leq a_\ell,b_\ell<p$, then
\[\binom ab\bmod p=\prod_{\ell=0}^w\binom{a_\ell}{b_\ell},\tag{4}\]
where we use the convention that $\binom cd=0$ when $d>c$. Thus, $\binom ab\bmod p\neq0$ iff $a_\ell\geq b_\ell$ for all $\ell=0,\ldots,w-1$.
Finally, we let $g\in\mathbb F_q[X]$ denote the polynomial $g(X)=X^q-X$. For $\boldsymbol X=(X_1,\ldots,X_m)$ and $\boldsymbol I=(i_1,\ldots,i_m)\in\mathbb N^m$ we let $g(\boldsymbol X)^{\boldsymbol I}$ denote $\prod_{k=1}^m(g(X_k))^{i_k}$.
\paragraph{Lemma 18.} Let $P\in\mathbb F[\boldsymbol X]$ be a multivariate polynomial where $\boldsymbol X=(X_1,\ldots,X_m)$. Let $\boldsymbol a,\boldsymbol b\in\mathbb F^m$ and define a bivariate polynomial by $P|_{P_{\boldsymbol a,\boldsymbol b,\boldsymbol c}}(t,r)=P(\boldsymbol at+\boldsymbol br+\boldsymbol c)$. Then for $(j_1,j_2)\in\mathbb N^2$:
\[(P|_{P_{\boldsymbol a,\boldsymbol b,\boldsymbol c}})^{(j_1,j_2)}(t,r)=\sum_{\boldsymbol I\in\mathbb N^m}P^{(\boldsymbol I)}(\boldsymbol at+\boldsymbol br+\boldsymbol c)\cdot\sum_{\substack{\boldsymbol I_1+\boldsymbol I_2=\boldsymbol I\\w(\boldsymbol I_1)=j_1,w(\boldsymbol I_2)=j_2}}\binom{\boldsymbol I}{\boldsymbol I_1}\boldsymbol a^{\boldsymbol I_1}\boldsymbol b^{\boldsymbol I_2}.\]
\paragraph{Proof.} Given $R_1,R_2\in\mathbb F$, we write the expression $P(\boldsymbol a(t+R_1)+\boldsymbol b(r+R_2)+\boldsymbol c)$ in two different ways. On the one hand,
\begin{align*}
P(\boldsymbol a(t+R_1)+\boldsymbol b(r+R_2)+\boldsymbol c)
&=P|_{P_{\boldsymbol a,\boldsymbol b,\boldsymbol c}}(t+R_1,r+R_2)\\
&=P|_{P_{\boldsymbol a,\boldsymbol b,\boldsymbol c}}((t,r)+(R_1,R_2))\\
&=\sum_{j_1,j_2\in\mathbb N}(P|_{P_{\boldsymbol a,\boldsymbol b,\boldsymbol c}})^{(j_1,j_2)}(t,r)R_1^{j_1}R_2^{j_2}.
\end{align*}
On the other hand,
\begin{align*}
P(\boldsymbol a(t+R_1)+\boldsymbol b(r+R_2)+\boldsymbol c)
&=P(\boldsymbol at+\boldsymbol br+\boldsymbol c+R_1\boldsymbol a+R_2\boldsymbol b)\\
&=\sum_{\boldsymbol I\in\mathbb N^m}P^{(\boldsymbol I)}(\boldsymbol at+\boldsymbol br+\boldsymbol c)\cdot(R_1\boldsymbol a+R_2\boldsymbol b)^{\boldsymbol I}\\
&=\sum_{\boldsymbol I\in\mathbb N^m}P^{(\boldsymbol I)}(\boldsymbol at+\boldsymbol br+\boldsymbol c)\cdot\prod_{k=1}^m(a_kR_1+b_kR_2)^{i_k}\\
&=\sum_{\boldsymbol I\in\mathbb N^m}P^{(\boldsymbol I)}(\boldsymbol at+\boldsymbol br+\boldsymbol c)
\sum_{\substack{\boldsymbol I_1+\boldsymbol I_2=\boldsymbol I\\w(\boldsymbol I_1)=j_1\\w(\boldsymbol I_2)=j_2}}\binom{\boldsymbol I}{\boldsymbol I_1}\boldsymbol a^{\boldsymbol I_1}\boldsymbol b^{\boldsymbol I_2}R_1^{j_1}R_2^{j_2}.
\end{align*}
Comparing coefficients of $R_1^{j_1}R_2^{j_2}$ for every $\boldsymbol J=(j_1,j_2)\in\mathbb N^2$ we get the result.\hfill$\triangleleft$
We will also be interested in the restrictions of polynomials to general $k$-dimensional subspaces. Let $P|_{P_{\boldsymbol h_0,\boldsymbol h_1,\ldots,\boldsymbol h_k}}(\boldsymbol Y)=P(\boldsymbol h_0+\sum_{i=1}^k\boldsymbol h_iY_i)$. Then, similarly to Lemma 18,
\paragraph{Lemma 19.}
\[(P|_{P_{\boldsymbol h_0,\boldsymbol h_1,\ldots,\boldsymbol h_k}})^{\boldsymbol J}(\boldsymbol Y)=\sum_{\boldsymbol I\in\mathbb N^m}P^{(\boldsymbol I)}\left(\boldsymbol h_0+\sum_{i=1}^k\boldsymbol h_iY_i\right)\cdot\sum_{\substack{\boldsymbol I_1+\cdots+\boldsymbol I_k=\boldsymbol I\\w(\boldsymbol I_r)=j_r}}\binom{\boldsymbol I}{\boldsymbol I_1,\ldots,\boldsymbol I_k}\prod_{i=1}^k\boldsymbol h_k^{\boldsymbol I_k}.\]
The proof is identical to the proof of Lemma 18.
\subsection*{2.3 Grobner bases and Nullstellensatz}
The theory of Grobner bases describes the structure of ideals in the ring $R=\mathbb F[\boldsymbol X]$ and we briefly explain some of the essential concepts of this theory. We refer to [5] for a thorough treatment of this theory.
\paragraph{Definition 20.} A monomial order $\succ$ on $R$ is a relation $\succ$ on $\mathbb Z_{\geq0}^n$, or equivalently a relation on the set of monomials $x^\alpha$, $\alpha\in\mathbb Z_{\geq0}^n$ satisfying:
\begin{enumerate}
\item $\succ$ is a total ordering.
\item If $\alpha\succ\beta$ and $\gamma\in\mathbb Z_{\geq0}^n$ then $\alpha+\gamma\succ\beta+\gamma$.
\item $\succ$ is a well-ordering. I.e, every non-empty $A\subset\mathbb Z_{\geq0}^n$ has a minimal element.
\end{enumerate}
\paragraph{Example 21 (Lexicographic order).} Let $\alpha,\beta\in\mathbb Z_{\geq0}^m$. We say $\alpha\succ_{\mathrm{lex}}\beta$ if the minimal $i$ which satisfies $\alpha_i\neq\beta_i$, also satisfies $\alpha_i>\beta_i$.
\paragraph{Example 22 (Total degree lexicographic order).} The total degree lexicographic order is defined as follows: A monomial $m_1$ is greater than $m_2$ if it has higher total degree, where ties are broken lexicographically (i.e $X_1>X_2>\cdots>X_m$). More formally, let $\alpha,\beta\in\mathbb Z_{\geq0}^m$. Then $\alpha\succ_{\mathrm{tot}}\beta$ if $w(\alpha)=\sum\alpha_i>w(\beta)=\sum\beta_i$ or, $w(\alpha)=w(\beta)$ and $\alpha\succ_{\mathrm{lex}}\beta$.
\paragraph{Definition 23.} Let $f(\boldsymbol X)=\sum_{\boldsymbol I}a_{\boldsymbol I}\boldsymbol X^{\boldsymbol I}$ and $\succ$ a monomial order.
\begin{enumerate}
\item The multi-degree of $f$ is $\operatorname{multideg}(f)=\max\{\boldsymbol I\mid a_{\boldsymbol I}\neq0\}$ where the maximum is taken w.r.t $\succ$.
\item The leading coefficient of $f$ is $\mathrm{LC}(f)=a_{\operatorname{multideg}(f)}\in\mathbb F$.
\item The leading monomial of $f$ is $\mathrm{LM}(f)=\boldsymbol X^{\operatorname{multideg}(f)}$.
\item The leading term of $f$ is $\mathrm{LT}(f)=\mathrm{LC}(f)\cdot\mathrm{LM}(f)$.
\end{enumerate}
\paragraph{Definition 24 (Multivariate polynomial division).} Let $\succ$ be a monomial order on $\mathbb Z_{\geq0}^m$, and let $F=\{f_1,\ldots,f_k\}$ be a set of $k$ polynomials in $\mathbb F[\boldsymbol X]$. Then every $f\in\mathbb F[\boldsymbol X]$ can be written as
\[f=q_1f_1+\cdots+q_kf_k+r,\]
where $q_i,r\in\mathbb F[\boldsymbol X]$, and either $r=0$ or $r$ is a linear combination, with coefficients in $\mathbb F$, of monomials, none of which is divisible by any of $\mathrm{LT}(f_1),\ldots,\mathrm{LT}(f_k)$. We call $r$ a remainder of the division by $F$. Moreover, $\operatorname{multideg}(q_if_i)\leq\operatorname{multideg}(f)$ for every $i\in[s]$. The remainder $r$ is not necessarily unique, and might depend on the order of division.
\paragraph{Definition 25.} Let $\{0\}\neq I\subseteq\mathbb F[\boldsymbol X]$ be an ideal. Fix a monomial ordering on $\mathbb F[\boldsymbol X]$. Then
\begin{enumerate}
\item We denote by $\mathrm{LT}(I)$ the set of leading terms of non-zero elements in $I$.
\[\mathrm{LT}(I)=\{\mathrm{LT}(f)\mid f\in I\setminus\{0\}\}.\]
\item We denote by $\langle\mathrm{LT}(I)\rangle$ the ideal generated by the elements in $\mathrm{LT}(I)$.
\end{enumerate}
\paragraph{Definition 26.} Let $\{0\}\neq I\subseteq\mathbb F[\boldsymbol X]$ be an ideal. Fix a monomial order on $\mathbb F[\boldsymbol X]$. A subset $G=\{g_1,\ldots,g_t\}\subset I$ is said to be a Grobner basis for $I$, if
\[\langle\mathrm{LT}(g_1),\ldots,\mathrm{LT}(g_t)\rangle=\langle\mathrm{LT}(I)\rangle.\]
\paragraph{Note 27.} Every ideal $I\subset\mathbb F[X_1,\ldots,X_m]$ is finitely generated and has a Grobner basis.
The importance of a Grobner basis, is that it gives a natural way of choosing representatives for the quotient space $\mathbb F[X_1,\ldots,X_m]/I$.
\paragraph{Theorem 28 ([5], Section 2, Proposition 1).} Let $I\subset\mathbb F[X_1,\ldots,X_m]$ be an ideal and $G=\{g_1,\ldots,g_k\}$ a Grobner basis. Then given $f\in\mathbb F[\boldsymbol X]$ there is a unique $r\in\mathbb F[\boldsymbol X]$ such that there is a $g\in I$ such that $f=g+r$ and no term of $r$ is divisible by any of $\mathrm{LT}(g_1),\ldots,\mathrm{LT}(g_k)$. We call this $r$, the reduced form of $f$ (relative to $I$).
Note that the reduced form of any polynomial is equivalent to this polynomial modulo $I$. Thus, the theorem gives us a natural way of choosing representatives modulo $I$.
\paragraph{Theorem 29.} Let $R=\mathbb F[\boldsymbol X]$ be the ring of polynomials, and $I\subset R$ an ideal. Let $G$ be a Grobner basis for $I$. Then the set
\[B=\{M(\boldsymbol X)\mid M\text{ is a monomial not divisible by any }\mathrm{LT}(g)\text{ for }g\in G\},\]
is a basis for $R/I$.
The following criterion determines whether $G$ is a Grobner basis.
\paragraph{Definition 30 (LCM and S polynomials).} Let $f,g\in\mathbb F[\boldsymbol X]$ be non-zero polynomials. Let $\alpha=\operatorname{multideg}(f)$ and $\beta=\operatorname{multideg}(g)$.
\begin{enumerate}
\item The least common multiple of $\mathrm{LM}(f)$ and $\mathrm{LM}(g)$, denoted $\mathrm{LCM}(\mathrm{LM}(f),\mathrm{LM}(g))$, is $\boldsymbol X^\gamma$, where $\gamma=(\gamma_1,\ldots,\gamma_m)$ and $\gamma_i=\max\{\alpha_i,\beta_i\}$ for each $i$.
\item The $S$-polynomial of $f$ and $g$ is
\[S(f,g)=\frac{\mathrm{LCM}(\mathrm{LM}(f),\mathrm{LM}(g))}{\mathrm{LT}(f)}\cdot f-\frac{\mathrm{LCM}(\mathrm{LM}(f),\mathrm{LM}(g))}{\mathrm{LT}(g)}\cdot g.\]
\end{enumerate}
\paragraph{Theorem 31 (Buchberger's Criterion, [5], Sec 6).} Let $I\subset\mathbb F[\boldsymbol X]$ be an ideal. Then a basis $G=\{g_1,\ldots,g_k\}$ of $I$ is a Grobner basis of $I$ if and only if for all pairs $i\neq j$, the remainder on division of $S(g_i,g_j)$ by $G$ is zero.
Note that we always have $S=S(g_i,g_j)\in I$ by the definition of $S$. When saying the remainder of the division by $G$ is zero, we mean that there are $\{f_i\}$, such that $S=\sum f_ig_i$ and $\operatorname{multideg}(f_ig_i)\leq\operatorname{multideg}(S)$ for every $i$ (as in Definition 24).
\paragraph{Theorem 32 (Combinatorial Nullstellensatz, [1]).} Let $\mathbb F$ be a field, and $A_1,\ldots,A_m\subseteq\mathbb F$. Let $g_i(X)=\prod_{\alpha\in A_i}(X-\alpha)$ for $i=1,\ldots,m$. Assume a polynomial $f\in\mathbb F[\boldsymbol X]$ satisfies $f(\alpha)=0$ for all $\alpha\in A_1\times\cdots\times A_m$. Then there are $h_1,\ldots,h_t$ such that
\[f=\sum h_ig_i,\]
and $\deg(h_i)+\deg(g_i)\leq\deg(f)$ for all $i$.
When $A_i=\mathbb F_q$ denote
\[g(X)=\prod_{\alpha\in\mathbb F_q}(X-\alpha)=X^q-X.\]
Also, let $I_m$ denote the ideal $I_m=\{f\in\mathbb F_q[\boldsymbol X]\mid\forall\alpha\in\mathbb F_q^m\ f(\alpha)=0\}$.
\paragraph{Corollary 33.} $I_m=\langle g(X_1),\ldots,g(X_m)\rangle$.
\paragraph{Proof.} Let $f\in I_m$. By Theorem 32, taking $S_i=\mathbb F_q$ for every $i$, we get that $f=\sum h_ig_i$ for some $\{h_i\}$ and so $f\in\langle g_k\rangle_{k\in[m]}$. The other inclusion is trivial since $g(X_i)=X_i^q-X_i$ vanishes on $\mathbb F_q^m$ for every $i$.\hfill$\triangleleft$
Using Theorem 31 it can be easily proved that:
\paragraph{Claim 34.} $G=\{g(X_k)\}_{k=1}^m$ is a Grobner basis for $I_m$ relative to the total degree lexicographic order.
Let $s\in\mathbb N$ and let $I_{m,s}$ denote the ideal
\[I_{m,s}=\{f\in\mathbb F_q[\boldsymbol X]\mid\forall\alpha\in\mathbb F_q^m\ \mathrm{Mult}(f;\alpha)\geq s\}.\]
In this notation, $I_{m,1}=I_m$ defined before. For every $\boldsymbol I=(I_1,\ldots,I_m)\in\mathbb N^m$ define
\[g(\boldsymbol X)^{\boldsymbol I}=\prod_{k=1}^mg(X_k)^{I_k}.\]
\paragraph{Theorem 35 (Combinatorial Nullstellensatz with multiplicity, [3], Sec 3).} $I_{m,s}=\langle g(\boldsymbol X)^{\boldsymbol I}\rangle_{w(\boldsymbol I)=s}$. Furthermore, the set $G_{m,s}=\{g(\boldsymbol X)^{\boldsymbol I}\}_{w(\boldsymbol I)=s}$ is a Grobner basis for $I_{m,s}$.
\paragraph{Proof.} To see that $I_{m,s}$ is indeed an ideal, fix $f\in I_{m,s}$ and $h\in\mathbb F[\boldsymbol X]$. Then for $r<s$: $(hf)^{(r)}=\sum_{i=0}^rf^{(i)}\cdot h^{(r-i)}=0$ and so $hf\in I_{m,s}$. Also, clearly, $g(\boldsymbol X)^{\boldsymbol I}\in I_{m,s}$ for every $\boldsymbol I$ with $w(\boldsymbol I)=s$. We need to show $G_{m,s}$ is a Grobner basis for $I_{m,s}$,
[3, Section 3] show $G_{m,s}$ generates $I_{m,s}$ and, furthermore, $f=\sum_{b:w(\boldsymbol I)=s}g(\boldsymbol X)^{\boldsymbol I}h_{\boldsymbol I}$ for some $h_{\boldsymbol I}$ with $\deg(h_{\boldsymbol I})\leq\deg(f)-s\deg(g)$. In particular, this is true for the $S$ polynomials in Theorem 31. I.e, every such $S$ polynomial can be expressed as $S=\sum g(\boldsymbol X)^{\boldsymbol I}h_{\boldsymbol I}$ where $\deg(g(\boldsymbol X)^{\boldsymbol I}h_{\boldsymbol I})\leq\deg(S)$. By Buchberger's criterion $g(\boldsymbol X)^{\boldsymbol I}$ is a Grobner basis.\hfill$\triangleleft$
Finally, we look at equality modulo $I_{m,1}$.
\paragraph{Definition 36.} For $n_1,n_2\in\mathbb N$ we say $n_1=_{\mathbb F_q}n_2$ iff $x^{n_1}=x^{n_2}\bmod I_{1,1}$. Equivalently, $n_1=_{\mathbb F_q}n_2$ iff
\begin{itemize}
\item $n_1=n_2$, or,
\item $\min(n_1,n_2)>0$ and $n_1=n_2\bmod(q-1)$.
\end{itemize}
\paragraph{Definition 37.} Let $\boldsymbol A,\boldsymbol B\in\mathbb N^m$. We say $\boldsymbol A=_{\mathbb F_q}\boldsymbol B$ for iff $A_k=_{\mathbb F_q}B_k$ for every $1\leq k\leq m$.
We also record:
\paragraph{Claim 38.} Let $a,b\in\mathbb N$. If $a=_{\mathbb F_q}b$ and $a<q$ then $a\leq b$.
\paragraph{Proof.} If $b<q$ then $a=_{\mathbb F_q}b$ implies $a=b$. Otherwise $a<q\leq b$.\hfill$\triangleleft$
\section*{3 Polynomials and tables (continued author context)}
Note that every polynomial determines its table $\mathrm{EVAL}(P)$. However, two different polynomials might have the same table if their difference is the zero table on $\mathbb F_q^m$. From Theorem 35,
\[\operatorname{Ker}(\mathrm{EVAL}_{m,s})=I_{m,s}=\langle g(\boldsymbol X)^{\boldsymbol I}\rangle_{w(\boldsymbol I)=s}.\]
Since $I_{m,s}$ does not contain any non-zero polynomial of total degree less than $sq$, we conclude that:
\paragraph{Corollary 40.} $\mathrm{EVAL}_{m,s}$ is injective on polynomials of total degree less than $sq$.
\paragraph{Claim 41.} $\mathrm{EVAL}_{m,s}$ is onto $\Sigma_{m,s}^{q^m}$.
\paragraph{Proof.} We will show that
\[\dim(\mathbb F_q[X_1,\ldots,X_m]/I_{m,s})\geq\dim((\Sigma_{m,s})^{q^m})=q^m\cdot\binom{m+s-1}{m},\]
where the dimension is over $\mathbb F_q$. This implies that the image of $\mathrm{EVAL}_{m,s}$ is everything, and the mapping is injective.
To see that consider the set
\[B_{m,s}=\left\{g(\boldsymbol X)^{\boldsymbol I}\cdot\prod_{k=1}^mX_k^{j_k}\mid w(\boldsymbol I)<s,0\leq j_k<q\right\}.\tag{5}\]
The elements in $B_{m,s}$ have different multi degree and therefore are independent. Also elements in $B_{m,s}$ are monomials of degree smaller than $sq$, and therefore are $I_{m,s}$ reduced. Thus,
\begin{align*}
\dim(\mathbb F_q[\boldsymbol X]/I_{m,s})&\geq\dim\operatorname{Span}(B_{m,s})\\
&=|\{(\boldsymbol I,\boldsymbol J)\in\mathbb N^m\times\mathbb N^m\mid w(\boldsymbol I)<s,0\leq j_k<q\}|\\
&=\binom{m+s-1}m\cdot q^m,
\end{align*}
as desired.\hfill$\triangleleft$
\paragraph{Corollary 42.} The set $B_{m,s}$ is a basis for $\mathbb F[\boldsymbol X]/I_{m,s}$.
\paragraph{Example 43.} When $m=1$ the set $\{g(x)^ix^j\}_{i\in\mathbb N,j<q}$ is a basis of $\mathbb F[\boldsymbol X]$. Thus, for every $\ell\in\mathbb N$ there are $\beta_{\ell,i,j}$ such that
\[x^\ell=\sum_{i_1,i_2:i_1q+i_2\leq\ell,i_2<q}\beta_{\ell,i_1,i_2}g(x)^{i_1}x^{i_2}.\]
By comparing coefficients of $x^\ell$ we see that $\beta_{\ell,\lfloor\ell/q\rfloor,\ell\bmod q}=1$.
\paragraph{Claim 44.} For every table $T:\mathbb F_q^m\to\Sigma_{m,s}$ there exists a degree $q(s-1)+(q-1)m$ polynomial such that $\mathrm{EVAL}(P)=T$.
\paragraph{Proof.} By Claim 41, $T$ may be written as $\mathrm{EVAL}(P)$, where $P$ is an $\mathbb F_q$-combination of basis elements in $B_{m,s}$. By the definition of $B_{m,s}$, any basis element is of the form $g(\boldsymbol X)^{\boldsymbol I}\cdot\prod_{k=1}^mX_k^{j_k}$ where $w(\boldsymbol I)<s$ and $0\leq j_k<q$. As such, $\deg(P)\leq\deg(g(\boldsymbol X)^{\boldsymbol I})+\deg(\prod_{k=1}^mX_k^{j_k})\leq q(s-1)+(q-1)m$.\hfill$\triangleleft$
We will later need:
\paragraph{Lemma 45.} Let $A\in\mathbb F_q[t]$ and $B\in\mathbb F_q[r]$. Then,
\[\deg_t(A(t)B(r)\bmod I_{2,s})\leq\deg_t(A).\]
\paragraph{Proof.} Express $A$ in the basis $\{g(t)^it^j\}_{i\in\mathbb N,j<q}$,
\[A(t)=\sum_{(i,j)\in\mathcal A}\alpha_{i,j}g(t)^it^j,\]
and similarly for $B$,
\[B(r)=\sum_{(k,\ell)\in\mathcal B}\beta_{k,\ell}g(r)^kr^\ell,\]
for some sets $\mathcal A,\mathcal B$, $\alpha_{i,j},\beta_{k,\ell}\in\mathbb F_q$. Then,
\[A(t)B(r)\bmod I_{2,s}=\sum_{(i,j)\in\mathcal A,(k,\ell)\in\mathcal B:i+k<s}\alpha_{i,j}\beta_{k,\ell}g(t)^ig(r)^kt^jr^\ell.\]
Suppose $\deg_t(A(t)B(r)\bmod I_{2,s})$ is achieved on the monomial of $(i,j)\in\mathcal A,(k,\ell)\in\mathcal B$. In particular, this degree in $t$ is already achieved in $A$ for the monomial $(i,j)$. Thus, $\deg_t(AB)\leq\deg_t(A)$.\hfill$\triangleleft$
Next, we would like to give a purely algebraic criteria, which states when exactly a table belongs to the code $\mathrm{MRM}(q,m,d,s)$.
\paragraph{Definition 46.} Let $T\in\Sigma_{m,s}^{q^m}$ be a table. By Corollary 40 and Claim 41 there is a unique element $P_T\in\mathbb F[\boldsymbol X]/I_{m,s}$ such that $\mathrm{EVAL}_{m,s}(P_T)=T$. We call $P_T$ the representing polynomial of $T$.
\paragraph{Lemma 47.} Assume $d<sq$. Let $T\in\Sigma_{m,s}^{q^m}$ be a table, and $P_T$ its representing polynomial. Then
\[T\in\mathrm{MRM}(q,m,d,s)\Longleftrightarrow\deg(P_T)\leq d.\]
\paragraph{Proof.} First, assume $\deg(P_T)\leq d$. Then, by definition, since $\mathrm{EVAL}_{m,s}(P_T)=T$, we have $T\in\mathrm{MRM}(q,m,d,s)$. For the other direction, assume $T\in\mathrm{MRM}(q,m,d,s)$. Then there is some $Q\in\mathbb F[\boldsymbol X]$ of total degree $\leq d$ such that $\mathrm{EVAL}_{m,s}(Q)=T$. As $\deg(Q)\leq d<sq$, $Q$ is $I_{m,s}$ reduced and by Corollary 40 $Q$ is the representing polynomial of $T$.\hfill$\triangleleft$
\section*{Imported RM tester and numerical transfer (author context)}
\paragraph{Theorem 2 ([12]).} Let $k=\lceil(d+1)/(q-q/p)\rceil$. Given $f:\mathbb F_q^m\to\mathbb F_q$ let:
\begin{itemize}
\item $\mathrm{REJ}_{k,d}(f)$ denotes the rejection probability of the test, namely, the probability over a random $k$ dimensional affine space, that $f$ restricted to the space cannot be explained by a degree $k$ polynomial, and,
\item $\delta(f,\mathrm{RM}(q,m,d))$ be the distance of $f$ from the code $\mathrm{RM}(q,m,d)$.
\end{itemize}
Then
\[\mathrm{REJ}_{k,d}(f)\geq\min\{\alpha_0\cdot\delta(f,\mathrm{RM}(q,m,d)),c_0\}.\]
where $\alpha_0=q^k/2$ and $c_0=1/(2(k+1)q^{k+1})$.
Combining the result with Theorem 2 we get:
\paragraph{Corollary 52.} Let $d,q,s<q$ be positive integers and let $t=t_{q,d}=\lceil(d+1)/(q-q/p)\rceil$ be the RM testing dimension. Then,
\[\mathrm{REJ}_{t,d}^{\mathrm{MRM}}(T)\geq\min\left\{\frac13\left(1-\frac{s-1}q\right)\delta(T,\mathrm{MRM}(q,m,d,s)),\frac1{2(t+1)q^{t+1}}\right\}.\]
\paragraph{Proof.} Theorem 2 from [12] ensures that
\[\mathrm{REJ}_{t,d}(f)\geq\min\left\{\frac{q^t}2\delta(f,\mathrm{RM}(q,m,d)),\frac1{2(t+1)q^{t+1}}\right\}.\]
Using Theorem 51 we see that
\[\alpha'\geq\left(1-\frac{s-1}q\right)\frac1{1+q^{d/(q-1)}\frac1\alpha}\geq\frac13\left(1-\frac{s-1}q\right).\]
using the fact that $d/(q-1)\leq t$ and $q^{d/(q-1)}/\alpha\leq2$.\hfill$\triangleleft$
In particular, for $d<(q-q/p)$ the line test is a local characterization, as $t_{q,d}=1$.
We note that the assumption $s\leq d$ is quite natural, as derivatives with order higher than the degree must be identically zero. In contrast, it is not clear whether the assumption $s<q$ is indeed required, and we leave it for future study.
